\chapter{Related Work}
\label{chapter:related-work}

This chapter reviews research relevant to URL-only phishing detection and
GAN-based URL generation. The review concentrates on the papers supplied for
the project and on the design decisions that affect the present work. The
studies are not ranked by their headline accuracy alone. Their datasets,
preprocessing, class balance, split rules, model inputs, and evaluation
procedures differ, so a numerical result from one paper is not automatically
comparable with a result from another.

\section{Review Scope and Comparison Method}

The reviewed literature is divided into three groups. The first group studies
URL structure and traditional machine-learning features. The second group uses
character-level neural networks for URL classification. The third group applies
GANs to simultaneous detection, synthetic-data augmentation, or adversarial
evasion. Review papers were used to confirm the broader classification of
phishing-detection methods, but the analysis gives more weight to studies that
describe their datasets, architecture, and experimental procedure
\citep{qasim2023review}.

Each experimental paper is considered using the following questions:

\begin{itemize}
\item What information is given to the model: raw URL characters, lexical
features, generated feature vectors, or another representation?
\item How many URLs are used, where do they come from, and how are duplicates
or repeated hosts handled?
\item Are training and testing separated randomly, by time, by identity, or by
source?
\item What are the generator and discriminator architectures, and what purpose
does each serve?
\item Are generated URLs evaluated for validity, uniqueness, diversity, and
overlap, or only through loss and classifier scores?
\item Are accuracy, precision, recall, F1-score, ROC-AUC, FPR, and FNR reported?
\item Does the experiment isolate the effect of synthetic augmentation from
changes in real data, class balance, thresholds, and test data?
\end{itemize}

These questions are important because a model can obtain a strong result for
the wrong reason. Repeated hosts may allow memorisation, URL schemes may reveal
the source, balanced tests may hide poor warning precision in ordinary browsing,
and thresholds chosen on the final test set can make performance appear better
than it would be on new data.

\section{URL Features and Traditional Classifiers}

Kan and Nguyen Thi studied webpage classification using the URL without
fetching the page. Their method segmented addresses and represented URL
components, component lengths, orthographic patterns, token sequences, and
precedence relations. The study was not specifically a phishing detector, but
it demonstrated that the location and order of URL tokens contain useful
information and that URL-only features can supplement other representations
\citep{kan2005fast}. This supports the use of hostname, path, query, length,
character, and token measurements in the lexical branch of the present hybrid
model.

Sahingoz et al. developed a phishing detector from URL data using natural
language processing features and several machine-learning algorithms. Their
study is important to this project because it treats phishing detection as a
URL-text classification problem and provides a traditional baseline for later
character-level work \citep{sahingoz2019machine}. However, later analysis of
the dataset found many repeated hosts, illustrating why deduplication and
identity separation can materially change reported performance
\citep{alalyan2020robust}.

Wang compared two feature-extraction procedures on 10,000 benign and 10,000
malicious URLs. One procedure used tokenisation and vectorisation, while the
other used selected lexical features. Logistic Regression, SVM, Random Forest,
Gradient Boosting, and Bagging were evaluated. Token/vector features combined
with ensemble methods produced the strongest results; Random Forest achieved
97.35\% accuracy and 98.75\% AUC in the reported condition
\citep{wang2022malicious}. The result supports RF and SVM baselines, but it
comes from that paper's own balanced dataset and feature space. It cannot be
used as a direct target for an external or low-phishing deployment test.

The traditional literature therefore provides two lessons. First, carefully
constructed lexical and token features can be highly informative. Second, the
same feature set may not transfer when data sources or collection periods
change. This motivates comparing lexical baselines against neural models on
exactly the same real rows in this project.

\section{Character-Level Deep Learning}

Al-Alyan and Al-Ahmadi proposed PUCNN, a URL-only character CNN, and created the
Massive URL Phishing Detection dataset. After removing duplicate URLs and
repeated hosts, their main dataset contained 1,140,599 legitimate and 1,167,201
phishing URLs. The CNN used a maximum length of 256 characters, an alphabet of
69 symbols, a character embedding, one-dimensional convolution, global max
pooling, and dense layers \citep{alalyan2020robust}.

PUCNN achieved 95.78\% accuracy with a random split. More importantly, the
authors also ordered the phishing data by collection time: older URLs were used
for training and newer URLs for testing. Accuracy fell to 88.54\%, a decrease of
about 7.5 percentage points. The paper also showed that a simple HTTP/HTTPS
scheme rule could reach approximately 74\% accuracy on its collected data, so
the scheme was removed to reduce source-related bias. These experiments are
particularly relevant because they show that random internal evaluation can
overstate performance when phishing patterns change over time.

The character-CNN study supports the main neural input used in this project,
but it also exposes two limitations that the present experiments must address.
A fixed maximum length truncates some information, and even a model trained on
more than two million URLs can experience substantial temporal distribution
shift. A larger dataset is helpful, but size alone does not guarantee
generalisation.

\section{GAN-Based Phishing Research}

GAN research in phishing has followed several different objectives. Some
studies use the discriminator as the final detector. Some add generated URLs to
classifier training. Others deliberately search for adversarial samples that
evade a black-box detector. These objectives should not be treated as the same
experiment.

\subsection{PWDGAN: Black-Box Evasion}

PWDGAN was designed to generate adversarial malicious examples that bypass
black-box phishing classifiers. It used 10,017 malicious websites from
PhishTank and 11,300 benign websites from Alexa. URLs were converted to 150
numeric features. The generator contained dense hidden layers, while the
discriminator learned to imitate the output of an external phishing detector
and provide feedback to the generator \citep{bac2021pwdgan}.

The study attacked SVM, Decision Tree, Random Forest, Logistic Regression, and
MLP detectors. After 500 epochs, the reported true-positive rate for most of
the attacked classifiers fell to approximately zero on adversarial examples.
This is evidence that a detector can be vulnerable to synthetic feature
perturbations. It is not evidence that the generated examples are valid,
human-readable phishing URLs, nor does it show that adding those examples to a
fixed classifier training set improves external real-URL detection. PWDGAN is
therefore most useful to this project as motivation for classifier-challenge
testing and adversarial hard-sample selection.

\subsection{Conditional GAN with Joint Detection}

Kamran et al. proposed a semi-supervised conditional GAN that simultaneously
generates URLs and classifies them. The generator combines an original URL,
class label, and noise vector. Its convolutional encoder--decoder produces a
synthetic URL representation. The multi-headed discriminator predicts both
real-versus-generated status and benign-versus-malicious class. Reconstruction,
adversarial, and classification losses are combined during training
\citep{kamran2022conditional}.

The study sampled 50,000 balanced URLs of at most 200 characters, using 40,000
for training and 10,000 for testing. The model was trained for 200 epochs. It
reported 95.52\% accuracy, 95.08\% precision, 96.00\% sensitivity, 95.54\%
F1-score, and 95.52\% AUC on real URL classification. On a balanced real and
adversarial evaluation, the discriminator achieved 87.45\% accuracy and detected
74.9\% of the adversarial examples.

The conditional design is useful because generation and classification share
class information, and reconstruction can keep generated output close to an
input URL. At the same time, high similarity may mean the generator makes local
variations rather than exploring broad new phishing modes. The study reports
structural similarity, but does not provide the complete uniqueness, exact
training-overlap, and cross-source tests required by the present project.

\subsection{PDGAN: LSTM Generator and CNN Discriminator}

PDGAN uses an LSTM generator and a nine-layer CNN discriminator. The generator
learns ordered URL characters and the discriminator extracts local character
patterns. It was evaluated using the large MUPD dataset, which after
preprocessing contained 1,167,201 phishing and 1,140,599 legitimate URLs. The
authors removed duplicate URLs and repeated hosts before splitting the data
\citep{alahmadi2022pdgan}.

The reported final detector achieved 97.58\% accuracy and 98.02\% precision.
Experiments added generated phishing sets of different sizes; the paper reports
its strongest condition after adding 50,000 synthetic URLs. The work is close
to the objective of this project because it combines a sequence generator,
character CNN, and defensive augmentation. It also acknowledges that some
generated URL details have incorrect semantic information. This reinforces the
need to inspect generated samples and to report quality separately from the
discriminator's classification metrics.

PDGAN primarily evaluates its own GAN discriminator and uses its own split.
It does not provide a paired, multi-seed comparison across RF, SVM, LSTM, CNN,
RNN, and a lexical--character hybrid on unchanged external tests. The present
work extends the experimental focus rather than attempting to reproduce its
headline result.

\subsection{WGAN-GP Augmentation}

Pham et al. applied WGAN-GP to character-level malicious URL generation and
then trained LSTM and GRU classifiers with and without generated phishing URLs
\citep{pham2021generated}. The study used malicious URL sources and two Common
Crawl legitimate sources. It evaluated total datasets of 150,000, 300,000,
600,000, and 1.5 million samples. The generator used one-dimensional
convolutions and residual blocks, while the critic used Wasserstein distance
and a gradient penalty.

The reported effect depended on both dataset size and classifier. At 150,000
samples, LSTM accuracy increased from 96.08\% to 96.47\% with augmentation. At
600,000 samples it decreased from 97.28\% to 97.16\%, and at 1.5 million it
decreased from 97.67\% to 97.57\%. GRU performance decreased slightly in every
reported augmented condition. The authors concluded that generated data was
more useful when real training data was limited.

This is important evidence against assuming that augmentation must always
improve a classifier. It supports optional real-data-size experiments, but also
shows why the project should not change dataset size merely to find a winning
GAN configuration.

\subsection{Comparison of Basic GAN, DCGAN, WGAN, and SeqGAN}

Pham, Pham, and Ta compared basic GAN, DCGAN, WGAN, and SeqGAN on a dataset
containing two million legitimate and two million phishing URLs. One million
phishing URLs were used for discriminator training, and a test set contained
one million phishing and one million legitimate URLs. A second experiment
trained GANs on balanced real data and added generated URLs to RF and LSTM
training \citep{pham2023evaluation}.

The discriminator results differed from generator-quality results. Basic GAN,
DCGAN, WGAN, and SeqGAN obtained reported discriminator accuracies of 99.98\%,
99.88\%, 98.99\%, and 85.05\%, respectively. However, only about half of the
basic-GAN, DCGAN, and WGAN outputs passed the top-level-domain structural check,
whereas SeqGAN reached 92.22\%. The paper's examples also show severe malformed
character output from some continuous GAN variants. This demonstrates that a
strong discriminator does not prove that its generator produces good URLs.

The downstream augmentation effects were small and mixed. RF accuracy was
90.46\% on original balanced data and ranged from 90.28\% to 90.85\% after the
different GAN additions. LSTM accuracy was 96.45\% on original data and ranged
from 96.16\% to 96.53\% after augmentation. SeqGAN produced the strongest
top-level-domain validity and the largest RF reduction, while basic GAN produced
the largest LSTM reduction. In that paper these reductions were interpreted as
evidence of classifier deception; for defensive augmentation, the meaning is
different and must be tested on real held-out data.

One methodological caution is that the paper describes searching thresholds
for the best accuracy on the test set. A final deployment study should instead
select thresholds on validation or calibration data and use the test set once.
The project follows the latter procedure.

\subsection{Generative Round Network}

Wanda et al. proposed the Generative Round Network, which adds stochastic
rounding to the generative components. It used 14,000 records with an 80/20
train--test split. The reported GRN result was 95\% accuracy, 97\% precision,
97\% recall, and 97\% F1-score, compared with 94\% accuracy for the paper's
standard GAN condition \citep{wanda2025deep}.

The stochastic mechanism is intended to keep generator and discriminator
training balanced. The study analyses generator and discriminator losses, but
loss balance is only indirect evidence of URL realism and diversity. The
dataset is relatively small compared with the million-scale URL studies, and
the authors identify temporal datasets and further architectures as future
work. Its results are therefore useful as evidence that alternative discrete
operations may stabilise training, not as a direct comparison with this
project's models.

\section{Comparison of the Reviewed Studies}

Table~\ref{tab:related-work-comparison} summarises the studies most closely
related to the implemented system. Reported metrics are reproduced to describe
each paper; they should not be read as a common leaderboard.

\small
\begin{longtable}{@{}>{\raggedright\arraybackslash}p{0.16\textwidth}>{\raggedright\arraybackslash}p{0.21\textwidth}>{\raggedright\arraybackslash}p{0.23\textwidth}>{\raggedright\arraybackslash}p{0.29\textwidth}@{}}
\caption{Comparison of closely related URL-detection and GAN studies.}
\label{tab:related-work-comparison}\\
\toprule
Study & Model and data & Main reported result & Relevance and limitation \\
\midrule
\endfirsthead
\multicolumn{4}{c}{\tablename\ \thetable{} -- continued}\\
\toprule
Study & Model and data & Main reported result & Relevance and limitation \\
\midrule
\endhead
\midrule
\multicolumn{4}{r}{Continued on next page}\\
\endfoot
\bottomrule
\endlastfoot

Al-Alyan and Al-Ahmadi \citep{alalyan2020robust} & Character CNN; about 2.31M deduplicated legitimate and phishing URLs. & 95.78\% random-split accuracy; 88.54\% with date split. & Strong evidence for URL characters and temporal shift; no GAN augmentation. \\

Wang \citep{wang2022malicious} & RF, SVM, LR, GB, and Bagging; 20K balanced URLs; lexical and token/vector features. & RF with token/vector features: 97.35\% accuracy and 98.75\% AUC. & Supports traditional baselines; smaller balanced evaluation. \\

Conditional GAN \citep{kamran2022conditional} & Convolutional conditional generator and multi-task discriminator; 50K balanced URLs. & 95.52\% real-URL accuracy; 87.45\% real/adversarial accuracy. & Joint generation and classification; limited cross-source and diversity evidence. \\

PDGAN \citep{alahmadi2022pdgan} & LSTM generator and CNN discriminator; about 2.31M URLs after preprocessing. & 97.58\% accuracy and 98.02\% precision. & Closest defensive GAN design; quality analysis and paired cross-model testing remain limited. \\

WGAN-GP augmentation \citep{pham2021generated} & Convolutional WGAN-GP plus LSTM and GRU; 150K to 1.5M conditions. & Small LSTM gain at 150K; mixed or negative effects at larger sizes and for GRU. & Shows that augmentation benefit depends on real-data size and classifier. \\

Four-GAN evaluation \citep{pham2023evaluation} & Basic GAN, DCGAN, WGAN, and SeqGAN; 4M real URLs. & Discriminator and generator rankings differed; augmentation changes were small and mixed. & Strong motivation for separate quality tests; test-based threshold search is a limitation. \\

PWDGAN \citep{bac2021pwdgan} & Dense feature-space generator and surrogate discriminator; 21,317 URLs. & Most attacked classifiers approached 0\% TPR after 500 epochs. & Demonstrates adversarial weakness; does not establish useful, valid URL-text augmentation. \\

GRN \citep{wanda2025deep} & GAN with stochastic rounding; 14K URLs. & 95\% accuracy and 97\% F1-score. & Explores training balance; smaller random split and limited direct quality evidence. \\
\end{longtable}
\normalsize

Several patterns emerge from the comparison. Large random-split datasets often
produce strong internal results, but the temporal CNN experiment shows that
performance can fall sharply on later phishing URLs. Generator and
discriminator quality are separate: the best discriminator may belong to a GAN
that produces malformed strings. Augmentation effects can also be classifier-
and data-size-dependent. Finally, accuracy is reported more consistently than
FPR, FNR, external precision, uniqueness, or overlap, although those measures
are central to browser use.

\section{Research Gap}

The reviewed work demonstrates that URL-only models and GAN-based approaches
are promising, but it leaves a gap between a high internal score and evidence
that synthetic phishing URLs improve a practical detector. The main missing
elements are:

\begin{itemize}
\item a single frozen real training dataset shared by every compared classifier;
\item paired real-only and GAN-augmented conditions that change only the added
synthetic phishing rows;
\item several fixed augmentation percentages retained even when they do not
improve performance;
\item repeated seeds and reporting of variation rather than one favourable run;
\item explicit generated-URL validity, uniqueness, novelty, overlap, structural
realism, diversity, and classifier-challenge measurements;
\item internal, identity-disjoint, source-diverse, temporal or recently
collected, low-phishing, and long-legitimate-URL evaluations;
\item complete reporting of precision, recall, F1-score, ROC-AUC, FPR, FNR, and
confusion counts; and
\item a browser decision policy whose threshold is selected without using the
final test labels.
\end{itemize}

The project addresses this gap using a controlled experiment rather than
assuming that every synthetic sample is beneficial. Real-only models establish
the reference performance. GANs are trained only on phishing URLs from the
training split. Generated candidates pass fixed quality and overlap checks,
then the same accepted pool is supplied to the classifiers being compared.
Validation and test sets remain real and unchanged. Results are interpreted as
evidence of benefit, no effect, or degradation according to the measured data.

The hybrid classifier adds a further comparison not covered consistently in
the supplied GAN papers: character CNN features and lexical measurements are
combined in one detector, but the hybrid receives the same real URL identities
as the ordinary baselines. Weighted sampling is treated as a separate training
condition rather than as additional data. This distinction prevents the hybrid
from gaining hidden source access.

\section{Chapter Summary}

Prior research establishes that URL structure, lexical features, RF and SVM
models, and character-level neural networks can all support phishing detection.
GAN studies show that LSTM, convolutional, conditional, Wasserstein, sequential,
and stochastic-rounding designs can contribute to generation or discrimination.
They also show that GAN output may be malformed, that a
strong discriminator does not guarantee a strong generator, and that
augmentation may help, have no effect, or reduce performance.

The main opportunity is therefore methodological: compare real-only and
GAN-augmented classifiers under identical data and evaluation conditions, while
measuring generated URL quality separately. Chapter 4 describes how that
controlled methodology and the complete system were designed.
